\begin{problem}{F}{Fixing Table}{X seconds}

\noindent
You are given two tables $A$ and $B$ of size $2 \times n$. 
There are $T$ queries:
\begin{itemize}
  \item \texttt{1 r x y}: Invert all bits in row $r$ and columns from $x$ to $y$ (inclusive) of table $A$. Here, $r$ is either \fonttt{0} or \fonttt{1}, and $1 \le x \le y \le n$.
  \item \texttt{2}: Output the minimum number of operations required to transform table $A$ into table $B$ using the following operation:
    \begin{itemize}
      \item choose a subtable of $A$ and invert all bits in that subtable (change \fonttt{0} to \fonttt{1} and \fonttt{1} to \fonttt{0}).
    \end{itemize}
\end{itemize}

\noindent
Note that the subtable is defined by choosing two rows and two columns, and it includes all bits in the rectangle formed by those rows and columns.

\InputFile

\noindent
The first line of input contains an integer $n$ ($1 \le n \le 10^5$).\\
The second line contains the first row of table $A_0$ ($|A_0| = n$).\\
The third line contains the second row of table $A_1$ ($|A_1| = n$).\\
The fourth line contains the first row of table $B_0$ ($|B_0| = n$).\\
The fifth line contains the second row of table $B_1$ ($|B_1| = n$).\\
The sixth line contains an integer $T$ ($1 \le T \le 10^5$).\\
The next $T$ lines contain the queries in the format described above.

\OutputFile

\noindent
Output the minimum number of operations required to transform $A$ into $B$.

\Examples

\begin{example}
\exmp{1}{
5
11010
00101
10101
01010
3
2
1 0 2 4
2
}{2
1
}

\exmp{2}{
4
1010
0101
0101
1010
3
2
1 1 1 3
2
}{1
0
}

\pagebreak

\end{problem}